\documentclass[twoside,11pt]{article}

\usepackage[preprint]{jmlr2e}

\usepackage{enumitem}
\usepackage{graphicx,color}
\usepackage{hyperref}
\usepackage{url}
\usepackage{booktabs}
\usepackage{amsfonts}
\usepackage{nicefrac}
\usepackage{microtype}
\usepackage{bbm}

\usepackage{amsmath, amsfonts, amssymb}
\newtheorem{assumption}{Assumption}
\usepackage{mathtools}
\usepackage{graphicx}
\usepackage{bm}

\usepackage{physics}

\usepackage{algorithm}
\usepackage{algpseudocode}

% Common shortcuts
\def\mbb#1{\mathbb{#1}}
\def\mfk#1{\mathfrak{#1}}

\def\bN{\mbb{N}}
\def \C{\mbb{C}}
\def \R{\mbb{R}}
\def\bQ{\mbb{Q}}
\def\bZ{\mbb{Z}}
\def \cph{\varphi}
\renewcommand{\th}{\theta}
\def \ve{\varepsilon}
\newcommand{\mg}[1]{\| #1 \|}

% Often helpful macros
\newcommand{\floor}[1]{\left\lfloor#1\right\rfloor}
\newcommand{\ceil}[1]{\left\lceil#1\right\rceil}
\renewcommand{\P}{\mathbb P\qty}
\newcommand{\E}{\mathbb{E}\qty}
\newcommand{\Et}{\mathbb{E}_{\mathcal F_{t-1}} \qty}
\newcommand{\Pt}{\mathbb{P}_{\mathcal F_{t-1}} \qty}
\newcommand{\Cov}{\mathrm{Cov}\qty}
\newcommand{\Var}{\mathrm{Var}\qty}

\newcommand{\lr}{\mathrm{lr}}
\renewcommand{\grad}{\nabla} 
\newcommand{\1}{\mathbbm{1}}

% Sets
\usepackage{braket}

\graphicspath{{/}}
\usepackage{float}

\renewcommand{\set}[1]{\left\{#1\right\}}
\newcommand{\given}{\;\middle|\;}

\title{EGGROLL Convergence}
\date{\today}
\author{\name Rohan Mukherjee \email rohan.mukherjee@st-hildas.ox.ac.uk \\
\addr Department of Computer Science \\
University of Oxford}

\begin{document}
    \maketitle
    \section{Problem Setup}
    \citet{sarkar2025evolution} introduce EGGROLL as a low-rank evolutionary strategies method for scalable black-box optimization of large neural networks, with major applications being LLM fine tuning and post-training. In this excerpt, we prove that (symmetric) EGGROLL converges to a first order stationary point. 

    \begin{assumption}[Properties of $f$]
        We suppose that $f: \R^{d \times m} \to \R$ satisfies the following properties:
            \begin{enumerate}
                \item $f$ is differentiable and lower bounded, i.e. $f^* = \min_x f(x) > -\infty$
                \item $f$ is $L$-gradient Lipschitz, i.e.
                \begin{align*}
                    \mg{\grad f(X) - \grad f(Y)}_F \leq L \mg{X-Y}_F \quad \forall X,Y \in \R^{d \times m}
                \end{align*}
            \end{enumerate} 
    \end{assumption}

    \begin{definition}
        A point $X \in \R^{d \times m}$ is an $\ve$-first order stationary point if 
        \begin{align*}
            \mg{\grad f(X)}_F \leq \ve
        \end{align*}
    \end{definition}

    In this excerpt we work with symmetric EGGROLL, following the low-rank perturbation framework of \citet{sarkar2025evolution}.
    \begin{definition}[Symmetric EGGROLL Update]
        We define a rank $r$ symmetric zeroth order estimator as follows:
        \begin{align*}
            g_\sigma^{r}(X) \coloneqq \frac 1r \sum_{i=1}^r \frac{f(X + \sigma P_{i}) - f(X - \sigma P_{i})}{2\sigma} P_{i}
        \end{align*}
        using $P_{i} = a_{i}b_{i}^\top$ where $a_i \overset{i.i.d.}{\sim} \mathcal N(0, I_d)$ and $b_i \overset{i.i.d.}{\sim} \mathcal N(0, I_m)$, and $\sigma > 0$ is the smoothing radius.
    \end{definition}

    \begin{algorithm}[H]
    \caption{Rank-\(r\) Symmetric EGGROLL}
    \label{alg:symmetric-eggroll}
    \begin{algorithmic}
        \Require Initial point \(X_0 \in \mathbb R^{d\times m}\), horizon \(T\), stepsize \(\eta>0\), smoothing radius \(\sigma>0\), rank \(r\in \mathbb N\)
        \For{\(t = 0,\ldots,T-1\)}
            \State Sample $\set{a_{t,i}}_{i=1}^r \sim \mathcal N(0,I)$ and $\set{b_{t,i}}_{i=1}^r \sim \mathcal N(0,I)$ to compute $g_\sigma^r(X_t)$
                
            \State Update $X_{t+1}
                \gets
                X_t-\eta g_\sigma^r(X_t)$
        \EndFor
        \State \Return \(X_T\)
    \end{algorithmic}
    \end{algorithm}

    We state an informal version of the main result.
    \begin{theorem}
        Consider running Algorithm~\ref{alg:symmetric-eggroll}. Suppose $0 < \delta \leq 1/e$. Suppose
        \begin{align*}
            \sigma = \tilde{O}\qty(\frac{\ve}{dm}), \quad \eta = \tilde{O}\qty(\frac{r}{dm}).
        \end{align*}
        Then in
        \begin{align*}
            T = \tilde{\Omega}\qty(\frac{dmL(f(X_0) - f^*)}{r \ve^2})
        \end{align*}
        iterations (with each iteration using $2r$ function evaluations), with probability at least $1-\delta$, at least $3/4$ of the iterates are $\ve$-first order stationary points. 
    \end{theorem}

    \section{Proof Outline}
    First, we prove a few concentration bounds which will aid us in bounding the error terms. The lion's share of work went into proving Theorem~\ref{error_concentration}, which is the main tool used to get the running time above. We then introduce a preliminary lemma on function decrease, Lemma~\ref{symm-eggroll-function-decrease}. Because the update could either decrease function value if it aligns sufficiently with the gradient, or it could increase it but by an amount of the same order as other error terms, we use Lemma~\ref{half-two-gradient} and Lemma~\ref{gradient-dominates} to handle these two cases. Armed with these lemmata, Lemma~\ref{function-change-for-large-tau} is the remaining work to be done, and we arrive at the formal verison of our main theorem, Theorem~\ref{thm:symm-eggroll-conv} after a careful choice of parameters.

    \section{Concentration Bounds}
    Here we prove a novel concentration bound about the EGGROLL error. Throughout this excerpt, all norms of matrices is the Frobenius norm
    \begin{align*}
        \mg{A}_F = \qty(\sum_{i,j} A_{i,j}^2)^{1/2}.
    \end{align*}

    \subsection{Novel Concentration Bound}
    Throughout this paper, we denote
    \begin{align*}
        \Et[X ; A] \coloneqq \E[X \mathbb I_A \given \mathcal F_{t-1}]
    \end{align*}

    We will need the following lemma from \citet{jin2019short}.

    \begin{lemma}
        Let random vectors $x_1, \ldots, x_n \in \R^d$, and corresponding filtrations $\mathcal F_t = \sigma(X_1, \ldots, X_t)$ satisfy that 
        \begin{align*}
            \Et[x_t] = 0, \quad \Pt(\mg{x_t} \geq s) \leq 2e^{-\frac{s^2}{2\sigma^2}} \quad \forall s \in \R
        \end{align*}
        Then there exists an absolute constant $c$ such for any fixed $\delta > 0, \theta > 0$, with probability at least $1-\delta$,
        \begin{align*}
            \left\|\sum_{i=1}^n X_i\right\| \leq c\cdot \theta \sum_{i=1}^n \sigma_i^2 + \frac{1}{\theta} \log(\frac{2d}{\delta}).
        \end{align*}
    \end{lemma}

    This will help us later.
    \begin{lemma}
        Let $X$ follow $\chi^2_d$ distribution and $0 < \delta < 1/e$. Then there exists $L_{\delta}$ such that for every $d$, with
        \begin{align*}
            A \coloneqq \set{dL_{\delta}^2 \leq X \leq 3d\log(C/\delta)}
        \end{align*}
        we have
        \begin{align*}
            \E[(X-d) ; A] = 0, \quad \P(A) \geq 1-\delta
        \end{align*}
        for some absolute constant $C$.
    \end{lemma}
    \begin{proof}
        Let $f_d(x)$ be the $\chi^2_d$ density,
        \begin{align*}
            f_d(x) = \frac{1}{2^{d/2} \Gamma(d/2)} x^{d/2 -1}e^{-x/2}, \quad x \geq 0.
        \end{align*}
        A simple calculation shows that
        \begin{align*}
            (x-d)f_d(x) = -2\dv{x} (xf_d(x))
        \end{align*}
        Define $U = 3d \log(1/\delta)$ and define the set
        \begin{align*}
            A = \set{L \leq X \leq U}
        \end{align*}
        We want
        \begin{align*}
            \E[(X-d) ; A] = \int_L^U (x-d)f_d(x)dx = -2Lf_d(L) - 2Uf_d(U) \overset{\mathrm{want}}{=} 0.
        \end{align*}
        Therefore, we need 
        \begin{align*}
            Lf_d(L) &= Uf_d(U) \\ \iff L^d e^{-L} &= U^d e^{-U} \\
            \iff U - L &= d\log(U/L)
        \end{align*}
        Now define $\ell = L/d$ and $u = U/d$. Then we want to solve
        \begin{align*}
            du - d\ell &= d\log(u/\ell) \\
            \iff \ell e^{-\ell} &= u e^{-u}
        \end{align*}
        The function $g(x) = xe^{-x}$ is increasing on $[0,1)$ and decreasing on $(1,\infty)$. Therefore, as long as $u > 1$, which follows as $\delta < 1/e$, there exists a unique $\ell < 1$ satisfying the above equation. Now notice, as $e^{-\ell} \geq e^{-1}$,
        \begin{align*}
            \ell e^{-1} \leq \ell e^{-\ell} = 3\log(1/\delta) \delta^3 \leq 3e^{-1} \delta^2
        \end{align*} 
        Now we have that
        \begin{align*}
            \P(A^c) = \P(X \geq U) + \P(X \leq L)
        \end{align*}
         Recall the Chernoff bound For $0 < \alpha < 1$,
        \begin{align*}
            \P(X \leq \alpha d) \leq (\alpha e^{1-\alpha})^{d/2} \leq (e \alpha)^{d/2}
        \end{align*}
        Then we have
        \begin{align*}
            \P(X \leq L) \leq \P(X \leq 3d \delta^2) \leq (3e \delta^2)^{d/2} = (\sqrt{3e} \delta)^d
        \end{align*}
        When $\sqrt{3e} \delta < 1$, this is $< \sqrt{3e} \delta$, otherwise ths bound holds vacuously. Therefore $\P(X \leq L) \leq \sqrt{3e} \delta$. On the other hand, using that for $t < 1/2$,
        \begin{align*}
            \E[e^{t\mg{Z}^2}] = (1-2t)^{-d/2},
        \end{align*}
        along with Markov's inequality, we get
        \begin{align*}
            \P(X \geq 3d \log(1/\delta)) = \P(\exp(X/3) \geq \exp(d \log(1/\delta))) \leq 3^{d/2} \delta^d \leq \sqrt{3} \delta
        \end{align*}
        And thus $\P(A) \geq 1-C\delta$ for the absolute constant $C = \sqrt{3} + \sqrt{3e}$. 
    \end{proof}

    Here is the novel concentration bound. The proof is inspired by Proposition 2 of \citet{ren2023escaping} but with many important differences.

    \begin{theorem}\label{error_concentration}
        Let $\set{\mathcal F_t}_{t \geq -1}$ be a filtration, and let $\set{a_t}_{t \geq 0}, \set{b_t}_{t \geq 0} $ be two sequences of independent random vectors such that $a_t \sim \mathcal N(0, I_d)$, $b_t \sim \mathcal N(0, I_m)$, $a_t, b_t \in \mathcal F_t$ and $a_t, b_t$ are independent of $\mathcal F_{t-1}$. Let $\set{V_t}_{t \geq 0}$ be a sequence of random matrices such that $V_t \in \mathcal F_{t-1}$. For each $\tau \geq 0$, let 
        \begin{align*}
            E_\tau = \sum_{t=0}^{\tau-1} \qty(V_t - \langle V_t, a_tb_t^\top \rangle a_t b_t^\top)
        \end{align*}
        Then, there exists some absolute constants $C > 0, c > 0$ such that for any $\tau \in \bZ^+$, and $0 < \delta < 1/e$, for any $\theta > 0$, with probability at least $1-\delta$ given $\mathcal F_{t-1}$, 
        \begin{align*}
            \left \|E_{\tau} \right \|_F \leq c \cdot \theta \sum_{t=0}^{\tau-1} dm \log^{3}\qty(\frac{C\tau m}{\delta}) \mg{V_t}_F^2 + \frac{1}{\theta} \log(\frac{Cm^2d \tau}{\delta}).
        \end{align*}
    \end{theorem}

    \begin{proof}
        Let $V_t = P\Sigma Q^\top = \sum_{i=1}^r \sigma_i p_i q_i^\top$ be the singular value decomposition of $V_t$. Extend $\set{q_1, \ldots, q_r}$ to an orthonormal basis $\set{q_1, \ldots, q_m}$ of $\R^m$. Define the events
        \begin{align*}
            A_t = \set{\sqrt{d} L_\delta < \mg{a_t} < \sqrt{3d \log(C/\delta)}}
        \end{align*}
        and
        \begin{align*}
            B_{t,i} = \set{L_\delta < |b_t^\top q_i| < \sqrt{3\log(C/\delta)}}, \quad B_t = \bigcap_{i=1}^m B_{t,i}.
        \end{align*}
        Let $R$ be an orthogonal matrix $RR^\top = R^\top R = I$, and notice that $Ra$ has the same distribution as $a$, and that $A_t$ is the same event when $\mg{a_t}$ is replaced with $\mg{Ra_t}$. Therefore,
        \begin{align*}
            R\Et[a_ta_t^\top ; A_t]R^\top = \Et[(Ra_t)(Ra_t)^\top ; A_t] = \Et[a_ta_t^\top ; A_t]
        \end{align*}
        Now, if $M$ is a symmetric matrix such that $RMR^\top = M$ for all $R$, it must be that $M$ is a multiple of the identity. Therefore,
        \begin{align*}
            \Et[a_ta_t^\top ; A_t] = \alpha I_d.
        \end{align*}
        We then get
        \begin{align*}
            \alpha d &= \Tr(\Et[a_ta_t^\top ; A_t]) = \Et[\mg{a_t}^2 ; A_t] \\&= \Et[(\mg{a_t}^2 - d) ; A_t] + d \P(A_t) = d\P(A_t)
        \end{align*}
        The most important fact for this entire proof, is to notice that the events $B_{t,i}$ are indepedent. This is because, since $q_i \perp q_j$ for $i \neq j$, the $b_t^\top q_i$ and $b_t^\top q_j$ are independent given $\mathcal F_{t-1}$. Thus, if $i \neq j$,
        \begin{align*}
            \Et[(Q^\top b_t)(Q^\top b_t)^\top_{i,j} ; B_t] &= \Et[(q_i^\top b_t) (q_j^\top b_t) ; B_t] 
            \\&= \Pt(\bigcap_{k \neq i,j} B_{t,k}) \cdot \Et[q_i^\top b_t ; B_{t,i}] \cdot \Et[q_j^\top b_t ; B_{t,j}] \\
            &= 0
        \end{align*}
        because if $Z$ is $\mathcal N(0,1)$ and $A \in \sigma(Z)$ is a symmetric event, we have $\E[Z ; A] = 0$. When $i=j$,
        \begin{align*}
            \Et[(q_i^\top b_t)^2 ; B_t] = \Pt(\bigcap_{k \neq i} B_{t,k}) \cdot \Et[(q_i^\top b_t)^2 ; B_{t,i}] = \Pt(B_t)
        \end{align*}
        by our judicious choice of $B_{t,i}$ (cf. the lemma before this). 

        As $Q \in \mathcal F_{t-1}$, we thus have
        \begin{align*}
            \Et[Q^\top b_tb_t^\top ; B_t] &= \Et[(Q^\top b_t)(Q^\top b_t)^\top ; B_t] \cdot Q^\top  
            \\&= \Pt(B)I \cdot Q^\top  = \Pt(B)Q^\top .
        \end{align*}

    Therefore,
    \begin{align*}
        \Et[V_t - a_ta_t^\top V_t b_tb_t^\top ; A_t, B_t] &= V_t\Pt(A_t \cap B_t) - \Et[a_ta_t^\top P \Sigma Q^\top b_tb_t^\top ; A_t, B_t]
    \end{align*}
    Now,
    \begin{align*}
    \Et[a_ta_t^\top P \Sigma Q^\top b_tb_t^\top ; A_t, B_t] &= \Et[a_ta_t^\top ; A_t] P \Sigma \cdot \Et[Q^\top b_tb_t^\top ; B_t] \\
    &= P\Sigma Q^\top = V_t
    \end{align*}

    Therefore, upon defining
    \begin{align*}
        Q_t = (V_t - a_ta_t^\top V_t b_tb_t^\top) \1_{A_t \cap B_t}
    \end{align*}
    we have
    \begin{align*}
        \Et[Q_t] = 0.
    \end{align*}

    The next part of the proof is showing that $Q_t$ is $\mathrm{nSG}(C\sqrt{mn} \mathrm{polylog}(dm/\delta) \mg{V}_F)$. First, notice that on $B_t$,
    \begin{align*}
        \mg{b_t}^2 = \sum_{i=1}^m |q_i^\top b_t| \leq 3m \log(C/\delta)
    \end{align*}
    and recall that $A_t$ was defined so that
    \begin{align*}
        \mg{a_t} \leq \sqrt{3d \log(C/\delta)}.
    \end{align*}
    So on $A_t \cap B_t$ we have
    \begin{align*}
        \mg{a_t a_t^\top V_t b_t b_t^\top}_F = |a_t^\top V b_t| \mg{a_t} \mg{b_t} \leq |a_t^\top V b_t| \sqrt{9md \log^2(C/\delta)}
    \end{align*}
    Our main goal now is to show that $a_t^\top V b_t$ is subgaussian with parameter $O( \mathrm{polylog}(1/\delta)\mg{V}_F)$. To simplify notation, we will drop the $t$s and conditioning, it won't be important for the argument. 
    Recall that for $Z \sim \mathcal N(0, \sigma^2)$, if $t\sigma^2 < \frac 12$,
    \begin{align*}
        \E[e^{t Z^2}] = (1-2t\sigma ^2)^{-1/2}
    \end{align*}
    Then given $b$, since $a^T V b$ has distribution $\mathcal N(0, \mg{Vb}^2)$, we see immediately that whenever $t \mg{Vb}^2 < \frac 14$,
    \begin{align*}
        \E[e^{t(a^T V b)^2} \given b] = (1-2t \mg{Vb}^2)^{-1/2} \leq e^{2t \mg{Vb}^2}
    \end{align*}
    where the latter bound comes from, if $x \leq \frac 14$,
    \begin{align*}
        \log((1-2x)^{-1/2}) = -\frac 12 \log(1-2x) \leq \frac 12 \cdot  \frac{2x}{1-2x} \leq 2x
    \end{align*}
    While it is now tempting to use $\mg{Vb} \leq \mg{V}_F \cdot \mg{b}$, this leads to a bad bound. Instead, observe that $Vb$ is $\mathcal N(0, VV^\top)$ distributed. 

    Now recall that for $Z \sim \mathcal N(0, A)$, if $t < 1/(2 \lambda_{\max}(A))$,
    \begin{align*}
        \E[e^{t \mg{Z}^2}] = \det(I-2tA)^{-1/2}
    \end{align*}
    Therefore, as long as $t < 1/(8\sigma_{\max}(V))$, we have that
    \begin{align*}
        \E[e^{2t\mg{Vb}^2}] &= \prod_{i=1}^r (1-4t \sigma_i^2)^{-1/2} = \exp(\sum_{i=1}^r -\frac 12 \log(1-4t\sigma_i^2)) \\
        &\leq \exp(4t \sum_{i=1}^r \sigma_r^2) = \exp(4t\mg{V}_F^2)
    \end{align*}
    We are now almost done. Observe, on the event $B$,
    \begin{align*}
        \mg{Vb}^2 = b^\top V^\top V b = \sum_{i=1}^r \sigma_i^2 b^\top q_i \leq 3\log(C/\delta) \sum_{i=1}^r \sigma_i^2 = 3\log(C/\delta) \mg{V}_F^2
    \end{align*}
    Now, $t\mg{Vb}^2 \leq t 3\mg{V}_F^2 \log(C/\delta)$, so choosing $t = 1/(24\mg{V}_F^2 \log(C/\delta))$, we will satisfy all the bounds above ($\sigma_{\max}(V) \leq \mg{V}_F^2)$, and get
    \begin{align*}
        \E[e^{t(a^\top V b)^2}] = \E[\E[e^{t (a^\top Vb)^2} \mid b]] \leq e^{4t \mg{V}_F^2} \leq \exp(\frac 1{6\log(C/\delta)}) \leq e^{1/6} \leq 2
    \end{align*}
    Therefore, $a_t^\top V_t b_t$ is $\mathrm{nSG}(\sqrt{24 \log(C/\delta)} \mg{V}_F)$. Combining the previous bounds on $\mg{a}$ and $\mg{b}$, we get
    \begin{align*}
        \Et[\exp(\frac{\mg{a_ta_t^\top V_t b_tb_t^\top \1_{A_t \cap B_t}}_F^2}{216 md \log^3(C/\delta)\mg{V}_F^2})] \leq 2
    \end{align*}
    Therefore, using $\mg{Q_t}_F^2 \leq \mg{V_t}^2 + \mg{a_ta_t^\top V_t b_tb_t^\top \1_{A_t \cap B_t}}$, for some absolute constant $c'$,
    \begin{align*}
        \Et[\exp(\frac{\mg{Q_t}_F^2}{c' md \log^3(C/\delta)\mg{V}_F^2})] \leq 2
    \end{align*}
    This proves that $Q_t$ is $\mathrm{nSG}(c'\log^{3/2}(C/\delta) \mg{V}_F)$. 

    Thus, by Lemma 1, for some absolute constant $c$, with probability at least $1-\delta$,
    \begin{align*}
        \left \|\sum_{t=0}^{\tau-1} Q_t \right\| \leq c \cdot \theta \sum_{t=0}^{\tau-1} md \log^{3}(C/\delta) \mg{V_t}_F^2 + \frac{1}{\theta} \log \frac{2mn}{\delta}.
    \end{align*}
    Now, on $\bigcap_{t=0}^{\tau-1} A_t \cap B_t$,
    \begin{align*}
        \sum_{t=0}^{\tau-1} Q_t = \sum_{t=0}^{\tau-1} (V_t - a_ta_t^\top V_t b_tb_t^\top)
    \end{align*}
    As $\Pt(B_{t,i}) \geq 1-\delta$, by union bound $\Pt(B_t) \geq 1-m \delta$. Further, recall that $\Pt(A_t) \geq 1-\delta$. Therefore,
    \begin{align*}
        \Pt(\bigcap_{t=0}^{\tau-1} A_t \cap B_t) \geq 1 - \tau(m+1)\delta
    \end{align*}
    so, with probability at least $1-\delta - \tau(m+1)\delta$, 
    \begin{align*}
        \left \| E_\tau \right \| \leq c \cdot \theta \sum_{t=0}^{\tau-1} md \log^{3}(C/\delta) \mg{V_t}_F^2 + \frac{1}{\theta} \log(\frac{2md}{\delta}).
    \end{align*}
    After rescaling $\delta$ to $\delta / (1 + \tau (m+1))$, we have arrived at the desired result.
    \end{proof}
    

    \subsection{Sub-Weibell Random Variables}
    \begin{definition}
        Random variable $X \in \R$ is $\mathrm{SW}(K,\alpha)$ for some $K,\alpha > 0$ if 
        \begin{align*}
            \P(|X| \geq s) \leq 2\exp(-(s/K)^{1/\alpha}) \quad \forall s \geq 0
        \end{align*}
    \end{definition}
    We will need the following lemmas.
    \begin{lemma}
        Suppose $X$ and $Y$ are $SW(K_X, \alpha)$ and $SW(K_Y, \alpha)$ respectively. Then $XY$ is $SW(C(K_X \cdot K_Y), 2\alpha)$ and $X+Y$ is $SW(C(K_X+K_Y), \alpha)$ for some absolute constant $C$.
    \end{lemma}
    \begin{lemma}
        Suppose $X_1, \ldots, X_n$ are identically distributed $SW(K', \alpha)$. Then for some absolute constant $c > 0$, for all $s \geq ncK'$, we have
        \begin{align*}
            \P(\qty|\sum_{i=1}^n X_i| \geq s) \leq \exp(-\qty(\frac{s}{ncK'})^{1/\alpha}).
        \end{align*}
    \end{lemma}
    We will need to frequently bound sums of $\mg{P_t}^{2k}$. 
    \begin{lemma}\label{noise-norm}
        Suppose $a_0, \ldots, a_{\tau-1} \sim \mathcal N(0, I_d)$, and $b_0, \ldots, b_{\tau-1} \sim \mathcal N(0, I_m)$. Then for any $k > 0$ there exists absolute constants $c,C > 0$ such that for any $0 < \delta < 1$, w.p. $ \geq 1 -\delta$,
        \begin{align*}
            \sum_{t=0}^{\tau-1} \mg{a_ib_i^\top}^{2k}_F \leq \tau Cc^{2k} d^k m^k (1+\log(1/\delta)^{2k}) 
        \end{align*}
        When $\delta < 1/e$, 
        \begin{align*}
            \sum_{t=0}^{\tau-1} \mg{a_ib_i^\top}^{2k}_F \leq 2\tau Cc^{2k} d^k m^k \log(1/\delta)^{2k}.
        \end{align*}
    \end{lemma}

    \begin{proof}
        $(a_i)_k^2$ and $(b_j)_l^2$ are $SW(1,1)$. Thus their product is $SW(c, 2)$. Thus $\mg{a_ib_i^T}_F^2 = \sum_{k,l} (a_i)_k^2 (b_j)_k^2$ is $SW(cmd, 2)$ for some absolute constant $c$. Thus, $\mg{a_ib_i^\top}_F^{2k}$ is $SW(c^km^kd^k, 2k)$. Therefore by the lemma, for $s \geq \tau Cc^km^kd^k$,
        \begin{align*}
            \P(\sum_{t=0}^{\tau-1} \mg{a_ib_i^\top}_F^{2k} \geq s) \leq \exp(-\qty(\frac{s}{\tau Cc^kd^km^k})^{1/2k})
        \end{align*}
        the lemma follows by choosing $s = (1+\log(1/\delta)^{2k}) \tau Cc^kd^km^k$.
    \end{proof}

    \section{Function Decrease}
    First, we prove the following crucial lemma.
    \begin{lemma}[Function decrease for rank $r$ symmetric EGGROLL]\label{symm-eggroll-function-decrease}
        Suppose at each time $t$, the algorithm performs the update step with rank $r \leq dm$,
        \begin{align*}
            X_{t+1} = X_t - \eta g_u^{r}(X_t)
        \end{align*}
        where
        \begin{align*}
            g_\sigma^{r}(X_t) \coloneqq \frac{1}{r} \sum_{i=1}^r \frac{f(X_t + \sigma P_{t,i}) - f(X_t - \sigma P_{t,i})}{2\sigma} P_{t,i}
        \end{align*}
        where $P_{t,i} = a_{t,i}b_{t,i}^\top$ where each $a_{t,i},b_{t,i}$ are drawn i.i.d. from $\mathcal N(0, I_d)$, and $\mathcal N(0, I_m)$ respectively, and $\sigma > 0$ is the smoothing radius.

        Then there exists absolute constants $c>0,C \geq 1$ such that, for any $T \in \bZ^+$ and $1 \leq \tau' \leq T$, and $0 < \delta < 1/e$, then upon defining $\mathcal H_{\tau}(\delta)$ to be the event on which the inequality 
        \begin{align*}
            f(X_{\tau}) - f(X_0) &\leq -\frac{\eta}{2} \sum_{t=0}^{\tau-1} \frac{1}{r} \sum_{i=1}^r |\langle \grad f(X_t), P_{t,i} \rangle|^2 + \tau C c^4 \eta \sigma^2 L^2(dm)^2 \log(\frac T\delta)^4 \\
        &\quad + \eta^2 \frac{cdm\chi^4}{r} \sum_{t=0}^{\tau-1} \mg{\grad f(X_t)}^2 + \eta^2 \tau C c^6  \sigma^2 L^3 (dm)^3 \log(\frac T \delta)^6
        \end{align*}
    holds, (where $\chi \coloneqq \log(Cm^2 drT/\delta)$), we have 
    \begin{align*}
        \P(\bigcap_{\tau=1}^{\tau'} \mathcal H_{\tau}(\delta)) \geq 1 - \frac{3\tau'}{T}\delta
    \end{align*}
    \end{lemma}
    \begin{proof}
    For each $t \in \set{-1, \ldots, \tau}$, let
    \begin{align*}
        \mathcal F_t = \sigma(X_0, \set{a_{0,i}}_{i=1}^r, \ldots, \set{a_{t,i}}_{i=1}^r, \set{b_{0,i}}_{i=1}^r, \ldots, \set{b_{t,i}}_{i=1}^r)        
    \end{align*}
    be the sigma-algebra containing the information available at time $t$, and as usual write $P_{t,i} = a_{t,i}b_{t,i}^\top$. 

    Observe that, by mean value theorem, there exists some $s_{t,i} \in [-1,1]$ such that 
    \begin{align*}
        g_\sigma(X_t) = \frac 1r \sum_{i=1}^r \frac{f(X_t+\sigma P_{t,i}) - f(X_t - \sigma P_{t,i})}{2 \sigma} P_{t,i} = \frac{1}{r} \sum_{i=1}^r \grad f(X_t+\sigma s_{t,i} P_{t,i})P_{t,i}
    \end{align*}
    Then we have that
    \begin{align*}
        X_{t+1}-X_t &= - \eta \cdot \frac 1r \sum_{i=1}^r(\langle \grad f(X_t), P_{t,i} \rangle P_{t,i} + \langle \grad f(X_t+\sigma s_{t,i} P_{t,i}) - \grad f(X_t), P_{t,i} \rangle P_{t,i})
    \end{align*}

    \begin{align*}
        f(X_{t+1}) - f(X_t) &\leq  \langle X_{t+1} - X_t, \grad f(X_t) \rangle + \frac L2 \mg{X_{t+1} - X_t}^2 \\
        &\leq - \eta \cdot \frac{1}{r} \sum_{i=1}^r |\langle \grad f(X_t), P_{t,i} \rangle|^2 
        \\& \quad - \eta \cdot \frac{1}{r} \sum_{i=1}^r \langle \grad f(X_t), P_{t,i} \rangle \langle \grad f(X_t + \sigma_{t,i} P_{t,i}) - \grad f(X_t), P_{t,i} \rangle \\
        &\quad + \frac {L\eta^2}{2} \left \|\frac 1r \sum_{i=1}^r(\langle \grad f(X_t), P_{t,i} \rangle P_{t,i} + \langle \grad f(X_t+\sigma s_{t,i} P_{t,i}) - \grad f(X_t), P_{t,i} \rangle P_{t,i}) \right \|^2
    \end{align*}

    Now notice, by using $ab \leq (a^2+b^2)/2$,
    \begin{align*}
        - \frac{\eta}{r} \sum_{i=1}^r \langle \grad f(X_t), P_{t,i} \rangle \langle \grad f(X_t + \sigma_{t,i} P_{t,i}) - \grad f(X_t), P_{t,i} \rangle \leq \frac{\eta}{2r} \sum_{i=1}^r (|\langle \grad f(X_t \rangle), P_{t,i}|^2 + \sigma^2 L^2 \mg{P_{t,i}}^4)
    \end{align*}
    Also, by $\mg{a+b}^2 \leq 2\mg{a}^2 + 2\mg{b}^2$,
    \begin{align*}
        &\left \|\frac 1r \sum_{i=1}^r(\langle \grad f(X_t), P_{t,i} \rangle P_{t,i} + \langle \grad f(X_t+\sigma s_{t,i} P_{t,i}) - \grad f(X_t), P_{t,i} \rangle P_{t,i}) \right \|^2
        \\&\leq 2\left\|\frac 1r \sum_{i=1}^r \langle \grad f(X_t), P_{t,i} \rangle P_{t,i} \right\|^2 + \frac{2}{r} \sum_{i=1}^r \left \| \langle \grad f(X_t+\sigma s_{t,i} P_{t,i}) - \grad f(X_t), P_{t,i} \rangle P_{t,i}) \right \|^2 \\
        &\leq 2\left\|\frac 1r \sum_{i=1}^r \langle \grad f(X_t), P_{t,i} \rangle P_{t,i} \right\|^2 + \frac{2\sigma^2 L^2}{r} \sum_{i=1}^r \mg{P_{t,i}}^6
    \end{align*}
    Therefore,
    \begin{align*}
        f(X_{t+1}) - f(X_t) &\leq -\frac{\eta}{2r} \sum_{i=1}^r |\langle \grad f(X_t), P_{t,i} \rangle|^2 + \frac{\eta \sigma^2 L^2}{2r} \sum_{i=1}^r \mg{P_{t,i}}^4 \\
        &\quad + L\eta^2 \left \| \frac{1}{r} \sum_{i=1}^r  \langle \grad f(X_t), P_{t,i} \rangle P_{t,i} \right\|^2 + \frac{\eta^2 L^3 \sigma^2}{r} \sum_{i=1}^r \mg{P_{t,i}}^6
    \end{align*}
    Now,
    \begin{align*}
        \left \| \frac{1}{r} \sum_{i=1}^r  \langle \grad f(X_t), P_{t,i} \rangle P_{t,i} \right\|^2 \leq 2\left \| \frac{1}{r} \sum_{i=1}^r  (\grad f(X_t) - \langle \grad f(X_t), P_{t,i} \rangle P_{t,i}) \right\|^2 + 2\mg{\grad f(X_t)}^2
    \end{align*}
    By our proposition, with probability $1-\delta$,
    \begin{align*}
        \left \|\sum_{i=1}^r  (\grad f(X_t) - \langle \grad f(X_t), P_{t,i} \rangle P_{t,i}) \right\| \leq c \cdot \theta dmr \log^3 \qty(\frac{C r m}{\delta}) \mg{\grad f(X_t)}^2 + \frac{1}{\theta} \log(\frac{Cm^2 dr}{\delta})
    \end{align*}
    As $\grad f(X_t)$ is fixed given $\mathcal F_{t-1}$, optimizing over $\theta$, this is at most
    \begin{align*}
        \sqrt{cdmr} \log(Cm^2dr/\delta)^2 \mg{\grad f(X_t)}.
    \end{align*}
    Therefore for some irrelevant absolute constant $c$,
    \begin{align}\label{error_equation}
        \left \| \frac{1}{r} \sum_{i=1}^r  \langle \grad f(X_t), P_{t,i} \rangle P_{t,i} \right\|^2 &\leq c \frac{dm}{r} \log(Cm^2dr/\delta)^4 \mg{\grad f(X_t)}^2
    \end{align}
    By union bound, this will hold jointly for all $0 \leq t \leq \tau'-1$ with probability at least $1-\tau' \delta$.
    
    Now by old lemma, for some absolute constants $C,c$ possibly different from before, with probability at least $1-\delta$,
    \begin{align*}
        \frac{1}{r} \sum_{t=0}^{\tau-1} \sum_{i=1}^r \mg{P_{t,i}}^{2k} \leq 2\tau Cc^{2k} (dm)^k \log(1/\delta)^{2k}.
    \end{align*}
    If $k$ is fixed, by union bound, this will hold jointly for all $1 \leq \tau \leq \tau'$ with probability $1-\tau' \delta$. Therefore, it will hold for both $k=2$ and $3$, for all $1 \leq \tau \leq \tau'$, with probability at least $1-2\tau'\delta$. 

    Therefore,
    \begin{align*}
        f(X_{\tau}) - f(X_0) &\leq -\frac{\eta}{2r} \sum_{t=0}^{\tau-1} \sum_{i=1}^r |\langle \grad f(X_t), P_{t,i} \rangle|^2 + \tau C c^4 \eta \sigma^2 L^2(dm)^2 \log(1/\delta)^4 \\
        &\quad + \eta^2 \frac{cdm}{r} \log(Cm^2 dr/\delta)^4 \sum_{t=0}^{\tau-1} \mg{\grad f(X_t)}^2 + \eta^2 \tau C c^6  \sigma^2 L^3 (dm)^3 \log(1/\delta)^6
    \end{align*}
    holds for all $1 \leq \tau \leq \tau'$ with probability at least $1-3 \tau \delta$.
    Equivalently,
    \begin{align*}
        \P(\bigcap_{\tau=1}^{\tau'} \mathcal H_\tau(\delta)) \geq 1-3\tau' \delta
    \end{align*}
    Similarly, we see that, as when $\tau$ is fixed we need only Equation~\ref{error_concentration} to hold jointly over $0 \leq t \leq \tau-1$, while the bound on the perturbation norms hold with probability at least $1-2\delta$, so,
    \begin{align*}
        \P(\mathcal H_\tau(\delta)) \geq 1 - (\tau+2)\delta
    \end{align*}
    After rescaling $\delta$ to $\delta/T$, we have arrived at the desired result.
    \end{proof}

    We will need the following Bernstein inequality for sub-exponential random variables \citep[Theorem~2.8.1]{vershynin2018high}.
    \begin{lemma}[Bernstein concentration inequality]
        Let $X_1, \ldots, X_n$ be independent, mean-zero, $\sigma$-subexponential random variables, i.e.
        \begin{align*}
            \E[\exp(\frac{|X_i|}{\sigma})] \leq 2 \quad \text{for all $i \in [n]$}.
        \end{align*}
        Then there exists an absolute constant $c$ such that for any $s \geq 0$,
        \begin{align*}
            \P(\sum_{i=1}^n X_i \geq s) \leq \exp(-c \min\set{\frac{s^2}{n\sigma^2}, \frac s \sigma}).
        \end{align*}
    \end{lemma}

    We then prove the following lemma.
    \begin{lemma}
        There exists an absolute constant $c_2 \geq 1$ such that, if
        \begin{align*}
            t_f(\delta) = \frac{c_2}{m}\log(\frac{T}{\delta}), \quad \delta > 0,
        \end{align*}
        and 
        \begin{align*}
            \mathcal B_{t_0}(\delta ; k) = \bigcup_{t=t_0}^{t_0+k-1} \set{\frac 1m \sum_{i=1}^m |a_{t,i}^\top \grad f(X_t) b_{t,i}|^2 \geq \frac 12 \mg{\grad f(X_t)}^2}
        \end{align*}
        we have
        \begin{align*}
            \P(\mathcal B_{t_0}(\delta; k)) \geq 1 - \frac{\delta}{T}
        \end{align*}
        for $0 < \delta < 1$, $t_0 \geq 0$ and $k \geq t_f(\delta)$. 
    \end{lemma}
    \begin{proof}
        By using the same technique from the proof of Theorem 3, we know that the set of random variables $\set{\mg{\grad f(X_t)}^2 - |a_{t,i}\grad f(X_t) b_{t,i}|^2}_{i=1}^m$ are independent, $c\mg{\grad f(X_t)}^2$ subexponential for some absolute constant $c > 0$ given $\mathcal F_{t-1}$. 
        Now define
        \begin{align*}
            E_t = \set{\frac 1m \sum_{i=1}^m |a_{t,i}^\top  \grad f(X_t) b_{t,i}|^2 < \frac 12 \mg{\grad f(X_t)}^2}
        \end{align*}
        Therefore,
        \begin{align*}
            \Pt(E_t) &= \Pt(\frac{1}{m} \sum_{i=1}^m |a_{t,i}^\top \grad f(X_t) b_{t,i}|^2 \geq \mg{\grad f(X_t)}^2) 
            \\&= \Pt(\sum_{i=1}^m \qty(\mg{\grad f(X_t)}^2 - |a_{t,i}^\top \grad f(X_t) b_{t,i}|^2) > \frac m2 \mg{\grad f(X_t)}^2) \\&\leq \exp(-c' \min\set{\frac{m}{4c^2}, \frac{m}{2c}}) \leq \exp(-c'' m)
        \end{align*}
        for some absolute constant $c''$, where the last line follows by using Bernstein concentration inequality.

        Now,
        \begin{align*}
            \P(\bigcap_{t=t_0}^{t_0+k-1} E_t) &= \P(\bigcap_{t=t_0}^{t_0+k-1} \set{\frac 1m \sum_{i=1}^m |a_{t,i}^\top \grad f(X_t) b_{t,i}|^2 < \frac 12 \mg{\grad f(X_t)}^2})  \\
            &= \E[\prod_{t=t_0}^{t_0+k-2} \1_{E_t} \mathbb E_{\mathcal F_{t_0+k-2}}[\1_{E_{t_0+k-1}}]] \leq \cdots \leq \exp(c'' mk)
        \end{align*}
        therefore by choosing $k \geq t_f(\delta)$ we get
        \begin{align*}
            \P(\bigcap_{t=t_0}^{t_0+k-1} \set{\frac{1}{m} \sum_{i=1}^m |a_{t,i} \grad f(X_t) b_{t,i}|^2 < \frac 12 \mg{\grad f(X_t)}^2}) \leq \frac{\delta}{T},
        \end{align*}
        which completes the proof.
    \end{proof}
    
    \begin{lemma}
        Define
        \begin{align*}
            \mathcal A_t(\delta) \coloneqq \set{\mg{\grad f(X_{t+1}) - \grad f(X_t)} \leq \frac{\mg{\grad f(X_t)}}{8t_f(\delta)} + \eta \sigma L^2 \cdot \frac 1r \sum_{i=1}^r \mg{P_{t,i}}^3}
        \end{align*}
        where $t_f(\delta)$ is defined above. Then there exists an absolute constant $c_2$ such that, whenever $\eta$ satisfies,
        \begin{align}
            \eta L c_2 \sqrt{\frac{dm}{r}} \log(\frac{CdmrT}{\delta})^2 \leq \frac{1}{8t_f(\delta)},
        \end{align}
        we have
        \begin{align*}
            \P(\mathcal A_t(\delta)) \geq 1 - \frac{\delta}{T}
        \end{align*}
        for any $0 < \delta \leq 1/e$ and $T \in \bZ^+$. 
    \end{lemma}
    \begin{proof}
        By $L$-lipschitz of $\grad f$,
        \begin{align*}
            \mg{\grad f(X_{t+1}) - \grad f(X_t)} &\leq L\mg{X_{t+1} - X_t} \\
            &\leq \eta L \left \|\frac 1r \sum_{i=1}^r a_{t,i}a_{t,i}^\top \grad f(X_t) b_{t,i}b_{t,i}^\top \right \| \\
            &\quad + \eta L \left \|\frac 1r \sum_{i=1}^r \langle \grad f(X_t + \sigma s_{t,i} P_{t,i}) - \grad f(X_t), P_{t,i} \rangle P_{t,i}\right \|
        \end{align*}
        By $L$-lipschitz of $\grad f(X_t)$ and the Cauchy-Schwarz inequality,
        \begin{align*}
            \left \|\frac 1r \sum_{i=1}^r \langle \grad f(X_t + \sigma s_{t,i} P_{t,i}) - \grad f(X_t), P_{t,i} \rangle P_{t,i}\right \| \leq L \sigma \cdot \frac{1}{r} \sum_{i=1}^r \mg{P_{t,i}}^3
        \end{align*}
        Now for some absolute constant $c_2 > 0$ by Theorem~\ref{error_concentration}, with probability at least $1-\delta/T$,
        \begin{align*}
            \left \|\frac 1r \sum_{i=1}^r a_{t,i}a_{t,i}^\top \grad f(X_t) b_{t,i}b_{t,i}^\top \right \| \leq c_2 \log(CdmrT/\delta)^2 \sqrt{\frac{dm}{r}} \mg{\grad f(X_t)}
        \end{align*} 
        plugging in our condition on $\eta$, we arrive at the statement of the lemma.
    \end{proof}
    
    This next very important lemma is the key that will let us lower bound 
    \begin{align*}
        \sum_{t=0}^{\tau-1} \frac 1r \sum_{i=1}^r |\langle \grad f(X_t), P_t \rangle|^2 \geq \tilde \Omega \qty(\sum_{t=0}^{\tau-1} \mg{\grad f(X_t)})
    \end{align*}
    in some sense.

    Denote
    \begin{align*}
        \mathcal E(t_1, t_2, \delta) = \bigcap_{t=t_1}^{t_1+t_2-1} \set{\mg{\grad f(X_t)} \geq 8t_f(\delta) \eta \sigma L^2 \cdot \frac 1r \sum_{i=1}^r \mg{P_{t,i}}^3}
    \end{align*}

    \begin{lemma}\label{half-two-gradient}
        Let $0 < \delta \leq 1/e$ and $T \in \bZ^+$ satisfying $T \geq 2t_f(\delta) + 1$. Consider any positive integer $t_f' \leq 2t_f(\delta)$, and any $t_0 \in \set{0, \ldots, T-1-t_f'}$. Suppose $\eta$ satisfies the condition above, then on
        \begin{align*}
            \mathcal E(t_0, t_f', \delta) \cap \qty(\bigcap_{t=t_0}^{t_0+t_f'-1} \mathcal A_t(\delta))
        \end{align*}
        we have
        \begin{align*}
            \frac 12 \mg{\grad f(X_0)} \leq \mg{\grad f(X_t)} \leq 2 \mg{\grad f(X_t)}
        \end{align*} 
        for all $t \in \set{t_0, \ldots, t_0+t_f'-1}$. 
    \end{lemma}
    \begin{proof}
        By plugging in the definition for $\mathcal E_t(\delta)$ into the definition of $\mathcal A_t(\delta)$, we see that on the prescribed event,
        \begin{align*}
            \mg{\grad f(X_{t+1})-\grad f(X_t)} \leq \frac{\mg{\grad f(X_{t})}}{4t_f(\delta)}
        \end{align*}
        therefore
        \begin{align*}
            \qty(1-\frac{1}{4t_f(\delta)})\mg{\grad f(X_t)} \leq \mg{\grad f(X_{t+1})} \leq \qty(1+\frac{1}{4t_f(\delta)}) \mg{\grad f(X_t)}
        \end{align*}
        which leads to
        \begin{align*}
            \qty(1 - \frac{1}{4t_f(\delta)})^{t-t_0} \mg{\grad f(X_0)} \leq \mg{\grad f(X_t)} \leq \qty(1+\frac{1}{4t_f(\delta)})^{t-t_0} \mg{\grad f(X_0)}.
        \end{align*}
        for all $t \in \set{t_0, \ldots, t_0+t_f'-1}$. Now using $\log(1+x) \leq x$ for $x \geq 0$ and $-\log(1-x) \leq x/(1-x)$ for $0 \leq x < 1$, we know that
        \begin{align*}
            \qty(1 - \frac{1}{4x})^{2x} \geq e^{-2/3} \geq \frac 12, \quad \qty(1+\frac{1}{4x})^{2x} \leq e^{1/2} \leq 2
        \end{align*}
        Now recalling that $t_f' \leq 2t_f(\delta)$, we are done.
    \end{proof}

    \begin{lemma}
        There exists an absolute constant $c > 0$ such that for any $t \in \bN$, the event
        \begin{align*}
            \mathcal G_t(\delta) \coloneqq \set{\frac 1r \sum_{i=1}^r \mg{P_{t,i}}^3 \leq c (dm)^{3/2} \log(T/\delta)^3}
        \end{align*}
        holds with probability at least $1-\delta/T$ for $0 < \delta \leq 1/e$. 
    \end{lemma}
    \begin{proof}
        Obvious from Lemma~\ref{noise-norm} and the Cauchy-Schwarz inequality.
    \end{proof}

    \begin{lemma}\label{gradient-dominates}
        Let $0 < \delta \leq 1/e$ and $T \in \bZ^+$ such that $T > 2t_f(\delta) + 1$. Consider any positive integer $t_f' \leq 2t_f(\delta)$ and $t_0 \in \set{0, \ldots, T-1-t_f'}$. Suppose $\eta$ satifies the condition above. Then on the event 
        \begin{align*}
            \mathcal E^c(t_0, t_f', \delta) \cap \qty(\bigcap_{t=t_0}^{t_0+t_f'-1} \mathcal A_t(\delta)) \cap \qty(\bigcap_{t=t_0}^{t_0+t_f-1} \mathcal G_t(\delta))
        \end{align*}
        we have
        \begin{align*}
            \mg{\grad f(X_t)} \leq 16 t_f(\delta) \eta \sigma L^2 c(dm)^{3/2} \log(T/\delta)^3
        \end{align*}
        for all $t \in \set{t_0, \ldots, t_0+t_f'-1}$. 
    \end{lemma}
    \begin{proof}
        Let $t'$ be the first time in $\set{t_0, \ldots, t_0+t_f-1}$ such that the gradient is dominated, meaning
        \begin{align*}
            \mg{\grad f(X_{t'})} \leq 8t_f(\delta) \eta L^2 \sigma \frac{1}{r} \sum_{i=1}^r \mg{P_{t,i}}^3.
        \end{align*}
        By definition of $\mathcal G_{t'}(\delta)$, we must have
        \begin{align*}
            \mg{\grad f(X_{t'})} \leq 8t_f(\delta) \eta L^2 \sigma c(dm)^{3/2} \log(T/\delta)^3
        \end{align*}
        For $j < t'$, using the idea from a previous lemma,
        \begin{align*}
            \mg{\grad f(X_j)} \leq 2 \mg{\grad f(X_{t'})} \leq 16 t_f(\delta)\eta L^2 \sigma c(dm)^{3/2} \log(T/\delta)^3
        \end{align*}
        On the other hand, for $t \in \set{t', \ldots, t_0+t_f'-1}$, by using definition of $\mathcal A_t(\delta)$ and $\mathcal G_t(\delta)$, we get
        \begin{align*}
            \mg{\grad f(X_{t+1})} &\leq \qty(1 + \frac{1}{8t_f(\delta)})\mg{\grad f(X_t)} + \eta \sigma L^2 c(dm)^{3/2} \log(T/\delta)^3
            \\&=\qty(1-\frac{1}{8t_f(\delta)})^{t+1-t'} \mg{\grad f(X_{t'})} 
            \\&\quad + \sum_{i=0}^{t-t'} \qty(1-\frac{1}{8t_f(\delta)})^{t-t'-i} \eta \sigma L^2 c(dm)^{3/2} \log(T/\delta)^3 \\
            &\leq \qty(1+\frac{1}{8t_f(\delta)})^{t_f'} \mg{\grad f(X_t)} 
            \\&\quad + 8t_f(\delta) \qty(\qty(1+\frac{1}{8t_f(\delta)})^{t_f'} - 1) \eta \sigma L^2 c(dm)^{3/2} \log(T/\delta)^3 \\
            &\leq e^{1/4} 8t_f(\delta) \eta \sigma L^2 c(dm)^{3/2} \log(T/\delta)^3 
            \\&\quad+ 8t_f(\delta)(e^{1/4} - 1) \eta \sigma L^2 c(dm)^{3/2} \log(T/\delta)^3
            \\&\leq 16 t_f(\delta) \eta \sigma L^2 c(dm)^{3/2} \log(T/\delta)^3
        \end{align*}
        Where we used $t_f' \leq 2t_f(\delta)$ and $(1+1/(8x))^{2x} \leq e^{1/4}$ when $x > 0$. This completes the proof. 
    \end{proof}

    We are now almost ready to prove the main theorem. But first, we need the following lemma.

    \begin{lemma}[Function change for large $\tau$]\label{function-change-for-large-tau}
        Let $0 < \delta \leq 1/e$ and let $\tau \geq t_f(\delta)$ be arbitrary. Split $\set{0, \ldots, \tau-1}$ into $K \coloneqq \tau/t_f(\delta)$ intervals $J_0, \ldots, J_{K-1}$ of size $t_f(\delta)$. Let $I_1$ denote the set of indices where the gradient always dominates the noise term in $J_k$, i.e.
        \begin{align*}
            I_1 \coloneqq \set{k \,:\, \mg{\grad f(X_t)} \geq 8t_f(\delta) \eta \sigma L^2 \frac 1r \sum_{i=1}^r \mg{P_{t,i}}^3  \text{ for all $t \in J_k$}}
        \end{align*}
        Suppose we choose $\eta$ satisfying
        \begin{align*}
            \eta \leq \frac{1}{8t_f(\delta)} \min\set{\sqrt{\frac{r}{dm}} \cdot \frac{1}{c_2 L \log(Cdmr\tau/\delta)^2}, \frac{r}{dm} \cdot \frac{1}{8c\chi^4}}
        \end{align*}

        Then on the event
        \begin{align*}
            \mathcal E_\tau(\delta) \coloneqq \mathcal H_\tau(\delta) \cap \qty(\bigcap_{t=0}^{\tau-1} \mathcal A_t(\delta)) \cap \qty(\bigcap_{t=0}^{\tau-1} \mathcal G_t(\delta) ) \cap \qty(\bigcap_{t=0}^{\tau-1} \mathcal B_{kt_f(\delta)}(\delta ; t_f(\delta)))
        \end{align*}
        we have the following upper bound for some absolute constants $c,C$,
        \begin{align*}
            f(X_\tau) - f(X_0) &\leq  -\sum_{k \in I_1} \frac{\eta}{4} \min_{t \in J_k} \mg{\grad f(X_t)}^2 + \tau \eta^4 \frac{cdm\chi^4}{r} t_f(\delta)^2 \sigma^2 L^4 (dm)^3 \log(T/\delta)^6\\
            &\quad + \tau C c^4 \eta \sigma^2 L^2(dm)^2 \log(\frac T\delta)^4 + \eta^2 \tau C c^6  \sigma^2 L^3 (dm)^3 \log(\frac T \delta)^6.
        \end{align*}
        Furthermore, 
        \begin{align*}
            \P(\mathcal E_\tau(\delta)) \geq 1 - \frac{(4\tau+2)}{T}\delta.
        \end{align*}
    \end{lemma}

    \begin{proof}
        Any interval $J_k = \set{t_0, \ldots, t_0+t_f-1}$ belongs to one of the following two cases:
        
        \textbf{(Case 1)} (Gradient dominates noise): This means for every $t \in J_k$,
        \begin{align*}
            \mg{\grad f(X_t)} \geq 8t_f(\delta) \eta \sigma L^2 \frac 1r \sum_{i=1}^r \mg{P_{t,i}}^3
        \end{align*}
        By the Lemma~\ref{half-two-gradient}
        we have
        \begin{align*}
            \min_{t \in J_k} \mg{\grad f(X_t)} \geq \frac 14 \max_{t \in J_k} \mg{\grad f(X_t)}
        \end{align*}
        On the event $\mathcal B_{kt_f(\delta)}(\delta ; t_f(\delta))$, there exists some $t \in J_k$ such that 
        \begin{align*}
            \frac 1r \sum_{i=1}^r |\langle P_{t,i}, \grad f(X_t) \rangle|^2 \geq \frac 12 \mg{\grad f(X_t)}^2
        \end{align*}
        Therefore, using an average,
        \begin{align*}
            \frac 14 \sum_{t \in J_k} \frac 1r \sum_{i=1}^r |\langle P_{t,i}, \grad f(X_t) \rangle|^2 &\geq \frac 14 \min_{t \in J_k} \mg{\grad f(X_t)} \geq \frac{1}{64} \max_{t \in J_k} \mg{\grad f(X_t)}^2
            \\& \geq \frac{1}{64t_f(\delta)} \sum_{t\in J_k} \mg{\grad f(X_t)}^2
        \end{align*}
        By choosing 
        \begin{align*}
            \frac{\eta cdm \chi^4}{r} \leq \frac{1}{64 t_f(\delta)}
        \end{align*}
        we have
        \begin{align*}
            &-\frac{\eta}{2} \sum_{t \in J_k} \frac 1r \sum_{i=1}^r |\langle P_{t,i}, \grad f(X_t) \rangle|^2 + \frac{\eta^2 cdm \chi^4}{r}\sum_{t \in J_k} \mg{\grad f(X_t)}^2 \\
            &\leq -\frac{\eta}{4} \sum_{t \in J_k} \frac 1r \sum_{i=1}^r |\langle P_{t,i}, \grad f(X_t) \rangle|^2 \leq -\frac{\eta}{4} \min_{t \in J_k} \mg{\grad f(X_t)}^2.
        \end{align*}

        \textbf{(Case 2)} (Gradient does not dominate noise): This means that there exists some $t \in J_k$ such that
        \begin{align*}
            \mg{\grad f(X_t)} \leq 8t_f(\delta) \eta \sigma L^2 \frac 1r \sum_{i=1}^r \mg{P_{t,i}}^3
        \end{align*}
        Then by Lemma~\ref{gradient-dominates}, we know that, for all $t \in J_k$,
        \begin{align*}
            \mg{\grad f(X_t)} \leq 16 t_f(\delta) \eta \sigma L^2 c(dm)^{3/2} \log(T/\delta)^3
        \end{align*}

        Every $J_k$ falls into one of the above two cases. Now we can complete the proof.
        \begin{align*}
            &-\frac{\eta}{2} \sum_{t=0}^{\tau-1} \frac 1r \sum_{i=1}^r |\langle P_{t,i}, \grad f(X_t) \rangle|^2 + \eta^2 \frac{cdm\chi^4}{r} \sum_{t=0}^{\tau-1} \mg{\grad f(X_t)}^2
            \\&= \sum_{k \in I_1} \qty(-\frac{\eta}{2} \sum_{t \in J_k} \frac 1r \sum_{i=1}^r |\langle P_{t,i}, \grad f(X_t) \rangle|^2 + \eta^2 \frac{cdm\chi^4}{r} \sum_{t \in J_k} \mg{\grad f(X_t)}^2)
            \\&\quad+\sum_{k \in I_1^c} \qty(-\frac{\eta}{2} \sum_{t \in J_k} \frac 1r \sum_{i=1}^r |\langle P_{t,i}, \grad f(X_t) \rangle|^2 + \eta^2 \frac{cdm\chi^4}{r} \sum_{t \in J_k} \mg{\grad f(X_t)}^2) \\
            &\leq -\sum_{k \in I_1} \frac{\eta}{4} \min_{t \in J_k} \mg{\grad f(X_t)}^2 + \sum_{k \in I_1^c} t_f(\delta) \qty(\eta^4 \frac{cdm\chi^4}{r} 256 t_f(\delta)^2 \sigma^2 L^4 c (dm)^3 \log(T/\delta)^6)
            \\&\leq -\sum_{k \in I_1} \frac{\eta}{4} \min_{t \in J_k} \mg{\grad f(X_t)}^2 + \tau  \eta^4 \frac{cdm\chi^4}{r} 256 t_f(\delta)^2 \sigma^2 L^4 c (dm)^3 \log(T/\delta)^6
        \end{align*}
        and thus by the definition of $\mathcal H_\tau(\delta)$ in Lemma~\ref{symm-eggroll-function-decrease},
        \begin{align*}
            f(X_\tau) - f(X_0) &\leq  -\sum_{k \in I_1} \frac{\eta}{4} \min_{t \in J_k} \mg{\grad f(X_t)}^2 + \tau \eta^4 \frac{cdm\chi^4}{r} 256 t_f(\delta)^2 \sigma^2 L^4 c (dm)^3 \log(T/\delta)^6\\
            &\quad + \tau C c^4 \eta \sigma^2 L^2(dm)^2 \log(\frac T\delta)^4 + \eta^2 \tau C c^6  \sigma^2 L^3 (dm)^3 \log(\frac T \delta)^6.
        \end{align*}
        Furthermore,
        \begin{align*}
            \P(\mathcal E_\tau(\delta)^c) &\leq \P(\mathcal H_\tau(\delta)^c) + \sum_{t=0}^{\tau-1} \P(\mathcal A_t(\delta)^c) + \sum_{t=0}^{\tau-1} \P(\mathcal G_t(\delta)^c) + \sum_{k=0}^{K-1} \P(\mathcal B_{kt_f(\delta)}(\delta; t_f(\delta))^c) 
            \\&\leq \frac{(\tau+2)}{T}\delta + \frac{\tau}{T}\delta + \frac{\tau}{T}\delta + \frac{K}{T}\delta \leq \frac{(4\tau+2)}{T}\delta
        \end{align*}
        which completes the proof.
    \end{proof}

    \begin{theorem}[Symmetric EGGROLL Convergence]\label{thm:symm-eggroll-conv}
        Let $0 < \delta < 1/e$ be arbitrary. Suppose we choose $\sigma, \eta$ and $T$ such that
        \begin{align*}
            \sigma^2 &\leq \min \set{\frac{\ve^2}{1024 C c^4 L^2 (dm)^2 \log(T/\delta)^4}, \frac{\ve^2}{8096 rL^3 (dm)^2 \log(T/\delta)^2}} \\
            \eta &\leq \frac{1}{8t_f(\delta)} \min\set{\sqrt{\frac{r}{dm}} \cdot \frac{1}{c L \log(Cdmr\tau/\delta)^2}, \frac{r}{dm} \cdot \frac{1}{8c\chi^4}}, \\
            T &\geq \max \set{4, \frac{512t_f(\delta)(f(X_0) - f^*)}{\eta \ve^2}}.
        \end{align*}
        Then with probability at least $1-5\delta$ there are at most $T/4$ iterations for which $\mg{\grad f(X_t)} \geq \ve$. 
    \end{theorem}

    \begin{proof}
        Without loss of generality assume that $t_f(\delta)$ divides $T$ and split $\set{0, \ldots, T-1}$ into $K=T/t_f(\delta)$ intervals $J_0, \ldots, J_{K-1}$. Let $I_1$ denote the set of indices $k$ such that for every $t \in J_k$,
        \begin{align*}
            \mg{\grad f(X_t)} \geq 8t_f(\delta) \eta \sigma L^2 \frac{1}{r} \sum_{i=1}^r \mg{P_{t,i}}^3.
        \end{align*}
        Denote
        \begin{align*}
            \mathcal E_T(\delta) \coloneqq \mathcal H_T(\delta) \cap \qty(\bigcap_{t=0}^{T-1} \mathcal A_t(\delta)) \cap \qty(\bigcap_{t=0}^{T-1} \mathcal G_t(\delta)) \cap \qty(\bigcap_{k=0}^{K-1} \mathcal B_{kt_f(\delta)}(\delta ; t_f(\delta)))
        \end{align*}
        For the rest of the proof assume we are working on $\mathcal E_T(\delta)$. By Lemma~\ref{function-change-for-large-tau}, and our choices of $\eta$ and $\delta$ in the statement of the theorem, we have
        \begin{align*}
            f(X_T) - f(X_0) &\leq  -\sum_{k \in I_1} \frac{\eta}{4} \min_{t \in J_k} \mg{\grad f(X_t)}^2 + T \eta^4 \frac{cdm\chi^4}{r} t_f(\delta)^2 \sigma^2 L^4 (dm)^3 \log(T/\delta)^6\\
            &\quad + T C c^4 \eta \sigma^2 L^2(dm)^2 \log(\frac T\delta)^4 + \eta^2 T C c^6  \sigma^2 L^3 (dm)^3 \log(\frac T \delta)^6
        \end{align*}
        Assume that there are more than $T/4$ iterations with $\mg{\grad f(X_t)} \geq \ve$. Let $I_\ve$ denote the set of indices $k$ where there is a $t \in J_k$ with $\mg{\grad f(X_t)} \geq \ve$. Thus by pigeonhole principle, $|I_\ve| \geq T/4t_f(\delta)$. By our choices of parameters, it can be shown that
        \begin{align*}
            16 t_f(\delta) \eta \sigma L^2 c(dm)^{3/2} \log(T/\delta)^3 < \ve
        \end{align*}
        While by Lemma~\ref{gradient-dominates}, if $k \in I_1^c$, 
        \begin{align*}
            \mg{\grad f(X_t)} \leq 16 t_f(\delta) \eta \sigma L^2 c(dm)^{3/2} \log(T/\delta)^3
        \end{align*}
        Therefore, $I_\ve \subset I_1$. Now by Lemma~\ref{half-two-gradient}, for any $k \in I_1$, 
        \begin{align*}
            \frac 12 \mg{\grad f(X_{kt_f(\delta)})} \leq \mg{\grad f(X_t)} \leq 2\mg{\grad f(X_{kt_f(\delta)})}, \quad \forall t \in J_k
        \end{align*}
        This implies that for any $k \in I_\ve$, $\min_{t \in J_k} \mg{\grad f(X_t)}^2 \geq \frac{1}{16} \ve^2$, and consequently,
        \begin{align*}
            -\sum_{k \in I_1} \frac{\eta}{4} \min_{t \in J_k} \mg{\grad f(X_t)}^2 \leq -\sum_{k \in I_\ve} \frac{\eta}{2} \cdot \frac{\ve^2}{16} \leq -\frac{T}{4t_f(\delta)} \cdot \frac{\eta}{4} \cdot \frac{\ve^2}{16} = -\frac{T \eta \ve^2}{256 t_f(\delta)}
        \end{align*}
        Hence,
        \begin{align*}
            f(X_T) - f(X_0) &\leq  -\frac{T \eta \ve^2}{256t_f(\delta)} + T \eta^4 \frac{cdm\chi^4}{r} t_f(\delta)^2 \sigma^2 L^4 (dm)^3 \log(T/\delta)^6\\
            &\quad + T C c^4 \eta \sigma^2 L^2(dm)^2 \log(\frac T\delta)^4 + \eta^2 T C c^6  \sigma^2 L^3 (dm)^3 \log(\frac T \delta)^6
        \end{align*}
        By our choice of parameters,
        \begin{align*}
            &T \eta^4 \frac{cdm\chi^4}{r} t_f(\delta)^2 \sigma^2 L^4 (dm)^3 \log(T/\delta)^6 + T C c^4 \eta \sigma^2 L^2(dm)^2 \log(\frac T\delta)^4 
            \\&\quad+ \eta^2 T C c^6  \sigma^2 L^3 (dm)^3 \log(\frac T \delta)^6 \leq \frac{T\eta \ve^2}{512 t_f(\delta)}
        \end{align*}
        We thus get
        \begin{align*}
            f(X_T) - f(X_0) \leq -\frac{T \eta \ve^2}{512 t_f(\delta)}
        \end{align*}
        Therefore, as long as
        \begin{align*}
            T \geq \frac{512t_f(\delta)(f(x_0) - f^*)}{\eta \ve^2}
        \end{align*}
        we would get $f(X_T) < f^*$, a contradiction. Thus, we may conclude that on the event $\mathcal E_T(\delta)$ there are at most $T/4$ iterations for which $\mg{\grad f(X_t)} \geq \ve$. 

        Finally,
        \begin{align*}
            \P(\mathcal E_T^c) &\leq \P(\mathcal H_T^c) + \sum_{t=0}^{T-1} \P(\mathcal A_t(\delta)) + \sum_{t=0}^{T-1} \P(\mathcal G_t(\delta)) + \sum_{k=0}^{K-1} \P(\mathcal B_{kt_f(\delta)}(\delta ; t_f(\delta))) \\
            &\leq \frac{(T+2)\delta}{T} + \delta + \delta + \frac{K\delta}{T} \leq 5\delta.
        \end{align*}
    \end{proof}
    \bibliography{references}
\end{document}
